%! Permutations and Transposes — Unit 2, Lesson 2.4 — Homework (Answer Key)
\documentclass[10pt]{article}
\usepackage{linalg-article}
\usepackage{linalg-key}   % pulls in -boxes; do NOT also load -boxes
\newcommand{\TallMath}[1]{$\displaystyle #1\rule[-1.4em]{0pt}{3.2em}$}
\newcommand{\xx}{\mathbf{x}}
\newcommand{\bb}{\mathbf{b}}
\newcommand{\cc}{\mathbf{c}}
\newcommand{\PP}{P}
\newcommand{\tp}{\mathsf{T}}

\begin{document}

\pageheader{Unit 2, Lesson 2.4}{Homework --- Answer Key}

\vspace{0.12in}

\begin{notesbox}{Practice}
Show your steps. Remember: a permutation $\PP$ swaps rows; when a zero pivot blocks $A=LU$, solve
$\PP A\xx=\PP\bb$ (swap the same entries of $\bb$). The transpose $A^{\tp}$ flips rows into columns, a
symmetric matrix has $A^{\tp}=A$, and $(AB)^{\tp}=B^{\tp}A^{\tp}$.

\begin{enumerate}[leftmargin=1.9em, itemsep=10pt]
  \item \textbf{Permutation matrix.} Let $A=\begin{bmatrix} 1 & 2 & 3 \\ 4 & 5 & 6 \\ 7 & 8 & 9 \end{bmatrix}$.
    \begin{enumerate}[label=(\alph*), leftmargin=1.8em, itemsep=4pt]
      \item Write the $3\times3$ permutation $\PP$ that swaps rows $1$ and $3$.
      \item Compute $\PP A$.
      \item Show that $\PP^{\tp}\PP=I$.
    \end{enumerate}
\begin{work}
  \PP &= \begin{bsmallmatrix} 0&0&1\\0&1&0\\1&0&0 \end{bsmallmatrix} \\
  \PP A &= \begin{bsmallmatrix} 7&8&9\\4&5&6\\1&2&3 \end{bsmallmatrix} \\
  \PP^{\tp}\PP &= \PP\PP = I \ \ \text{(swapping twice)}
\end{work}
  \item \textbf{Fix the zero pivot and solve.} For each, write $\PP$, compute $\PP A$ and $\PP\bb$, then
        solve $A\xx=\bb$.
    \begin{enumerate}[label=(\alph*), leftmargin=1.8em, itemsep=4pt]
      \item $A=\begin{bmatrix} 0 & 3 \\ 1 & 4 \end{bmatrix}$, \; $\bb=\begin{bmatrix} 6 \\ 9 \end{bmatrix}$
      \item $A=\begin{bmatrix} 0 & 2 \\ 3 & 5 \end{bmatrix}$, \; $\bb=\begin{bmatrix} 2 \\ 8 \end{bmatrix}$
    \end{enumerate}
\begin{work}
  \text{(a)}\ \PP A = \begin{bsmallmatrix} 1&4\\0&3 \end{bsmallmatrix},\ \PP\bb &= \begin{bsmallmatrix} 9\\6 \end{bsmallmatrix},\quad (x,y) = (1,2) \\
  \text{(b)}\ \PP A = \begin{bsmallmatrix} 3&5\\0&2 \end{bsmallmatrix},\ \PP\bb &= \begin{bsmallmatrix} 8\\2 \end{bsmallmatrix},\quad (x,y) = (1,1)
\end{work}
  \item \textbf{A $3\times3$ with a zero pivot.} Let $A=\begin{bmatrix} 0 & 1 & 2 \\ 2 & 4 & 6 \\ 2 & 5 & 11 \end{bmatrix}$
        and $\bb=\begin{bmatrix} 3 \\ 12 \\ 18 \end{bmatrix}$. Write $\PP$ to fix the pivot, factor
        $\PP A=LU$, then solve in two sweeps ($L\cc=\PP\bb$, then $U\xx=\cc$).
\begin{work}
  \PP A &= \begin{bsmallmatrix} 2&4&6\\0&1&2\\2&5&11 \end{bsmallmatrix},\quad
      \PP\bb = \begin{bsmallmatrix} 12\\3\\18 \end{bsmallmatrix} \\
  U &= \begin{bsmallmatrix} 2&4&6\\0&1&2\\0&0&3 \end{bsmallmatrix},\quad
      L = \begin{bsmallmatrix} 1&0&0\\0&1&0\\1&1&1 \end{bsmallmatrix} \\
  \cc &= \begin{bsmallmatrix} 12\\3\\3 \end{bsmallmatrix} \\
  (x,y,z) &= (1,1,1)
\end{work}
  \item \textbf{Transpose and symmetry.}
    \begin{enumerate}[label=(\alph*), leftmargin=1.8em, itemsep=4pt]
      \item Write $A^{\tp}$ for $A=\begin{bmatrix} 1 & 2 & 3 \\ 4 & 5 & 6 \end{bmatrix}$.
      \item Which of these are symmetric? \;
            $B=\begin{bmatrix} 2 & 0 \\ 0 & 7 \end{bmatrix},\quad
             C=\begin{bmatrix} 3 & 5 \\ 5 & 3 \end{bmatrix},\quad
             D=\begin{bmatrix} 1 & 2 \\ 3 & 4 \end{bmatrix}.$
    \end{enumerate}
\begin{work}
  \text{(a)}\ A^{\tp} &= \begin{bsmallmatrix} 1&4\\2&5\\3&6 \end{bsmallmatrix} \\
  \text{(b)}\ B,\,C \text{ symmetric};&\ D \text{ is not } (2 \ne 3)
\end{work}
  \item \textbf{The reversal rule.} For $A=\begin{bmatrix} 1 & 0 \\ 2 & 1 \end{bmatrix}$ and
        $B=\begin{bmatrix} 1 & 4 \\ 0 & 1 \end{bmatrix}$, compute $AB$ and confirm
        $(AB)^{\tp}=B^{\tp}A^{\tp}$. (Also check that $A^{\tp}B^{\tp}$ gives a \emph{different} matrix.)
\begin{work}
  AB &= \begin{bsmallmatrix} 1&4\\2&9 \end{bsmallmatrix},\quad
      (AB)^{\tp} = \begin{bsmallmatrix} 1&2\\4&9 \end{bsmallmatrix} \\
  B^{\tp}A^{\tp} &= \begin{bsmallmatrix} 1&2\\4&9 \end{bsmallmatrix} \ \ \text{\checkmark} \\
  A^{\tp}B^{\tp} &= \begin{bsmallmatrix} 9&2\\4&1 \end{bsmallmatrix} \ \ \text{(different)}
\end{work}
  \item \textbf{Explain.} Why is a permutation matrix's transpose equal to its inverse, i.e.\
        $\PP^{\tp}=\PP^{-1}$? Answer in terms of what a swap does and how to undo it.
        \ansline{$\PP$ reorders rows; the transpose reverses that reordering (row $i\to j$ becomes $j\to i$),
        which is exactly the swap that undoes it. So $\PP^{\tp}\PP=I$, meaning $\PP^{\tp}=\PP^{-1}$. (For a single swap, $\PP^{\tp}=\PP$, and swapping twice is $I$.)}
\end{enumerate}
\end{notesbox}

\begin{extensionbox}
\textbf{$A^{\tp}A$ is always symmetric.} Take any $A=\begin{bmatrix} 1 & 2 \\ 0 & 1 \\ 1 & 0 \end{bmatrix}$.
Compute $S=A^{\tp}A$ and check that $S^{\tp}=S$. Then explain the general reason using the rules
$(XY)^{\tp}=Y^{\tp}X^{\tp}$ and $(A^{\tp})^{\tp}=A$: what does $(A^{\tp}A)^{\tp}$ simplify to?
\begin{work}
  A^{\tp} &= \begin{bsmallmatrix} 1&0&1\\2&1&0 \end{bsmallmatrix} \\
  S = A^{\tp}A &= \begin{bsmallmatrix} 2&2\\2&5 \end{bsmallmatrix} = S^{\tp} \\
  (A^{\tp}A)^{\tp} &= A^{\tp}(A^{\tp})^{\tp} = A^{\tp}A
\end{work}
\ansline{The transpose reverses the order and undoes the inner transpose, giving back $A^{\tp}A$ --- so it is symmetric.}
\end{extensionbox}

\begin{spiralbox}
\textbf{Unit 2 wraps up here.} You can now solve $A\xx=\bb$ start to finish --- elimination, $A=LU$, and
a permutation $\PP$ for the zero-pivot case ($\PP A=LU$). Next comes the \textbf{Unit 2 Test}, then
\textbf{Unit 3: The Four Fundamental Subspaces}, where we stop asking ``\emph{what} is the solution?'' and
start asking ``\emph{which} $\bb$ are reachable, and what do \emph{all} the solutions look like?''
\end{spiralbox}

\end{document}
